%! Solving Quadratics by Factoring — Unit 2, Lesson 2.2 — Lesson Plan
% SLIDES-FIRST plan (course shape 2026-09-25; 2.2 is the pilot): the deck carries the lesson —
% warm-up, four cycles of definition → worked example → Now you try, a wrap-up, and the homework
% (a printed DeltaMath set) started in class. Students take notes on the printed slide handout.
% Teacher-facing; never handed out.
\documentclass[10pt]{article}
\usepackage{algebra2-article}
\usepackage{algebra2-boxes}
\usepackage{graphicx}   % -article does NOT load graphicx; needed for thumbnails/screenshots
\graphicspath{{images/}}
\newcommand{\TallMath}[1]{$\displaystyle #1\rule[-1.4em]{0pt}{3.2em}$}
\newcommand{\CourseName}{Algebra 2}
\newcommand{\MeetingLength}{60 minutes}
\newcommand{\UnitNumberName}{Unit 2: Quadratic Functions \quad}
\newcommand{\LessonNumberName}{Lesson 2.2: Solving Quadratics by Factoring}

\begin{document}
\begin{center}
    {\Large\bfseries \CourseName \\
    {\small\bfseries \UnitNumberName \LessonNumberName}}
\end{center}

\begin{tcolorbox}[colback=forestbg, colframe=forest, boxrule=0.9pt, arc=2mm,
    left=2mm, right=2mm, top=1.5mm, bottom=1.5mm]
    \textbf{Primary Objective:} Students \emph{produce} the factored form of a quadratic --- GCF
    first, $x^2+bx+c$ by the product--sum search, $ax^2+bx+c$ by grouping, and the
    difference-of-squares and perfect-square patterns --- and use the Zero Product Property to
    solve $ax^2+bx+c=0$, reading the roots as the graph's $x$-intercepts. Lesson~2.1
    \emph{read} zeros off a given factored form; 2.2 \emph{finds} that form.
    \\[2pt]
    \textbf{Standards (2023 VA SOL):} \textbf{A2.EO.3b} (factor completely over the integers),
    \textbf{A2.EO.3d} (difference of squares, perfect-square trinomials), \textbf{A2.EI.2b}
    (solve a quadratic equation algebraically --- here by factoring; real, factorable cases
    only), \textbf{A2.EI.2d} (verify solutions graphically: the roots are the $x$-intercepts).
    \\[2pt]
    \textbf{Lesson model:} \emph{Slides first.} The deck is the lesson and its printed handout is
    the student's notes. A warm-up on the slides; then, for each of four ideas, a
    \textbf{definition}, a \textbf{worked example} (I do), and a \textbf{Now you try} worked
    alone in the handout's notes column before the answer is revealed on a click; a short
    wrap-up; and the homework started in class. \textbf{Homework source:} DeltaMath, printed ---
    the exported set is dropped into the packet as \texttt{homework/main.pdf}.
\end{tcolorbox}

\renewcommand{\arraystretch}{1.4}
\begin{skillbox}[Priority Ideas \& Skills]{goldbox}
\begin{tabularx}{\linewidth}{@{}X|X@{}}
\textbf{Priority Ideas \& Skills} & \textbf{Key Understandings (the ``why'')} \\[2pt]
\hline
\vspace{-6pt}
\begin{itemize}[leftmargin=1.1em, nosep, after=\vspace{2pt}]
  \item Take out the \textbf{GCF} first, then factor what remains completely.
  \item Factor $x^2+bx+c$: two integers with \textbf{product $c$} and \textbf{sum $b$}.
  \item Factor $ax^2+bx+c$ by \textbf{grouping} (product $a\cdot c$, sum $b$).
  \item Factor $a^2-b^2$ and $a^2\pm2ab+b^2$ on sight; know $a^2+b^2$ does not factor.
  \item Solve by the \textbf{Zero Product Property}: $=0$, factor, each factor $=0$.
\end{itemize}
&
\vspace{-6pt}
\begin{itemize}[leftmargin=1.1em, nosep, after=\vspace{2pt}]
  \item Factoring \emph{un-multiplies}: $(x+p)(x+q)=x^2+(p+q)x+pq$ is why the search works.
  \item A product is zero \emph{only} when a factor is zero --- the step that turns one
        quadratic into two linear equations.
  \item \textbf{Target misconception: the property is about zero.} Students divide both sides
        of $x^2=5x$ by $x$ (losing $x=0$), or set each factor of $(x-1)(x+2)=4$ equal to $4$.
        Both skip ``get it equal to zero first.''
  \item The roots are the $x$-intercepts --- the algebra lands on the 2.0--2.1 picture.
\end{itemize}
\end{tabularx}
\end{skillbox}

\begin{skillbox}[{Vocabulary, Concepts, \& Theorems --- one definition frame each}]{sky}
\renewcommand{\arraystretch}{1.3}
\begin{tabularx}{\linewidth}{@{}l X@{}}
\textbf{Factor} & Write a polynomial as a product. Take out the \textbf{greatest common factor
                  (GCF)} first. \\
\textbf{Product--sum search} & $x^2+bx+c=(x+p)(x+q)$ where $pq=c$ and $p+q=b$. \\
\textbf{Grouping} & For $ax^2+bx+c$: two integers with product $ac$ and sum $b$ split $bx$; group
                    in pairs; take out the common binomial. \\
\textbf{Difference of squares} & $a^2-b^2=(a+b)(a-b)$. (A \emph{sum} of squares does not
                                 factor over the integers.) \\
\textbf{Perfect-square trinomial} & $a^2+2ab+b^2=(a+b)^2$, \; $a^2-2ab+b^2=(a-b)^2$; its graph
                                    touches the axis once --- a \textbf{double root}. \\
\textbf{Zero Product Property} & If $AB=0$, then $A=0$ or $B=0$. \\
\textbf{Roots / zeros} & The solutions of $ax^2+bx+c=0$; the $x$-intercepts of
                         $y=ax^2+bx+c$. \\
\end{tabularx}
\end{skillbox}

\begin{fixedskillbox}[Lesson at a Glance --- \MeetingLength]{forestbg}
\renewcommand{\arraystretch}{1.25}
\begin{tabularx}{\linewidth}{@{}l c X X@{}}
\textbf{Phase} & \textbf{Min} & \textbf{Students} & \textbf{Teacher} \\
\hline
Warm-Up & 5 & Multiply $(x+1)(x-3)$; a product--sum pair; read two zeros --- in the notes column & Reveal the answers; land ``how do we get the product back?'' \\
Lesson & 40 & Four rounds: read the definition, watch the example, then \emph{Now you try} alone before the reveal & Four cycles of $\sim$10 min: factoring $x^2+bx+c$, grouping, special patterns, Zero Product Property; circulate during each try \\
Wrap-up & 5 & Read back the four ideas; the caution & The wrap-up frame; say the ``about zero'' caution aloud \\
Homework & 10 & Start the DeltaMath pages alone & Launch problem~1, then circulate \\
\end{tabularx}
\end{fixedskillbox}

\begin{skillbox}[Warm-Up (5 min, on the slides)]{forestbg}
Three items, answers revealed one click each. \textbf{(1)} Multiply $(x+1)(x-3)=x^2-2x-3$ ---
the move today undoes. \textbf{(2)} Product $-12$, sum $1$: $4$ and $-3$ --- a dry run of the
search. \textbf{(3)} $y=(x+1)(x-3)$ crosses at $-1$ and $3$ (2.1). Close on the \emph{Hold on to
this} block: given only $x^2-2x-3$, how do you get the product back? That is cycle~1.
\end{skillbox}

\begin{skillbox}[The Lesson --- definition, example, now you try (40 min)]{forestbg}
\begin{multicols}{2}
\textbf{Cycle 1 --- factoring $x^2+bx+c$, GCF first (10 min).} Definition: factoring is
un-multiplying; the identity $(x+p)(x+q)=x^2+(p+q)x+pq$ is \emph{why} product $c$ / sum $b$ works
--- point at it. Example: $x^2-2x-3$ with the pair table (the warm-up's product comes back), then
$2x^2-4x-6=2(x+1)(x-3)$. Now you try: \textbf{(1)} $x^2+5x-14=(x+7)(x-2)$; \textbf{(2)}
$3x^2+3x-18=3(x+3)(x-2)$. Give 3 minutes. Look for students who skip the GCF in (2) and stall on
a product of $-18$ --- it factors that way too, $(3x+9)(x-2)$, but not \emph{completely}.

\textbf{Cycle 2 --- grouping (10 min).} Definition: the three steps. Example: $2x^2+5x-3$,
$ac=-6$, pair $6,-1$, $2x(x+3)-1(x+3)=(2x-1)(x+3)$, checked by multiplying. Now you try:
$3x^2-10x+8$, pair $-6,-4$, $3x(x-2)-4(x-2)=(3x-4)(x-2)$. The slip is the second group's sign:
students take out $+4$ and get $(x-2)$ and $(-x+2)$ --- ``the brackets must match; if they are
opposites, change the sign of the GCF.''

\columnbreak
\textbf{Cycle 3 --- special patterns (8 min).} Definition: $a^2-b^2$ and $a^2\pm2ab+b^2$, with the
tangent parabola $(x-3)^2$ --- a double root, the graph touches rather than crosses. Example:
$x^2-25$, $4x^2-9$, $x^2-6x+9$. Now you try: $9x^2-16=(3x+4)(3x-4)$; $x^2+10x+25=(x+5)^2$;
$x^2+16$ \emph{does not factor}. Expect $(x+4)(x-4)$ or $(x+4)^2$ for the third --- multiply either
out to kill it.

\textbf{Cycle 4 --- the Zero Product Property: the crux (12 min).} Definition: $AB=0\Rightarrow
A=0$ or $B=0$; the three-step solve; the graph of $x^2-2x-3$ with its zeros. Example:
$x^2=2x+3\to x^2-2x-3=0\to x=-1,\,3$, checked, and matched to the graph. Now you try:
\textbf{(1)} $x^2=5x$ --- \emph{watch for division by $x$}; the answer is $x=0$ or $x=5$.
\textbf{(2)} $(x-1)(x+2)=4$ --- \emph{watch for ``$x=5$ or $x=2$''}; the answer is $x=-3$ or
$x=2$. Give 4 minutes; before revealing, ask a student with $x=5$ only and one with $x=0,5$ to
say how they got there. This is the one frame not to rush.
\end{multicols}
\end{skillbox}

\boxguard
\begin{skillbox}[Wrap-up (5 min)]{forestbg}
Read the four bullets of the wrap-up frame back with the class, then say the caution aloud:
\emph{the property is about zero --- get everything on one side equal to $0$ before you factor,
and never divide by $x$.} Ask for a quick thumbs: ``$(x-2)(x+5)=6$ --- can I set $x-2=6$?''
(No.)
\end{skillbox}

\boxguard
\begin{skillbox}[Homework --- scored, started in class, due the first class after two study halls]{goldbox}
\textbf{Source: DeltaMath, printed.} Export the set as a PDF into
\texttt{unit02/lesson02/homework/main.pdf} (and its key into \texttt{homework\_key/main.pdf}),
then rebuild the packet. Build the set to cover, in order: GCF + $x^2+bx+c$; $ax^2+bx+c$;
difference of squares and perfect squares; solving by factoring, including at least one
$x^2=kx$ and one equation not already set to zero. \textbf{Set name:} TODO. \textbf{Formative
check:} use the $x^2=kx$ problem --- sort into \emph{both roots} / \emph{lost $x=0$ (divided by
$x$)} / \emph{other}; if the middle pile is large, open 2.3 by re-solving it. \textbf{Preview:}
2.3 --- when a quadratic will not factor, square roots and completing the square.
\end{skillbox}

\boxguard
\begin{skillbox}[Active Monitoring --- Watch For]{redbox}
\begin{itemize}[leftmargin=1.2em, nosep]
  \item \textbf{NYT 1.2:} factoring $3x^2+3x-18$ without the GCF --- not \emph{completely}
        factored.
  \item \textbf{NYT 2:} the second group's sign in grouping; brackets that do not match.
  \item \textbf{NYT 3.3:} $x^2+16=(x+4)(x-4)$ --- multiply it out.
  \item \textbf{NYT 4.1:} dividing by $x$ and reporting only $x=5$ (\emph{the target
        misconception}).
  \item \textbf{NYT 4.2:} setting each factor equal to $4$.
\end{itemize}
Cold-call prompts: ``How do you know that is factored \emph{completely}?'' \; ``What number does
the Zero Product Property need on the right?'' \; ``Where is that root on the graph?''
\end{skillbox}

% ── Teacher notes — one per component, in the order they are used ──
\begin{teachernote}[Slides]
Pacing: warm-up 5, cycles 10 / 10 / 8 / 12, wrap-up 5, homework 10. Each cycle is definition
(1 min, read aloud), example (3--4 min, click through the steps), Now you try (3--4 min, alone,
in the notes column), reveal (1 min). Cycle 4's Now you try is the one that matters; if you are
behind, cut the third item of cycle 3 ($x^2+16$) and the check line of Example~2, never cycle 4.
Students write in the handout's notes column, which prints without the answers.
\end{teachernote}

\begin{teachernote}[Homework]
The homework is a printed DeltaMath set --- the lesson builds without it, but the packet is
cover-only until \texttt{homework/main.pdf} is in place. Launch problem~1 aloud in the last 10
minutes, then circulate with the four Watch For items in mind. The $x^2=kx$ problem is the
formative check (sort above). Due the first class after two study halls; announce the date in
class and on TurtleNet.
\end{teachernote}

\end{document}
